\ifx\allfiles\undefined
%生成目录超链接，使用UTF8字符集
\documentclass[12pt,UTF8]{ctexart}


\usepackage{mathtools}%花括号括起的公式
\usepackage{enumerate}%列表
\usepackage{graphicx}%表格
\usepackage[table]{xcolor}%彩色表格
\usepackage{booktabs}%调整表格行高


\renewcommand\tabcolsep{50pt}
\renewcommand\arraystretch{2}

%文本框
\usepackage{tcolorbox}
\tcbuselibrary{skins,vignette,raster,listings,listingsutf8,theorems, breakable,magazine,fitting}

%页面设置
\usepackage{geometry}%设置页面
\usepackage{setspace}%设置局部行距

\usepackage{ctex}
\usepackage{CJK}

%图像处理
\usepackage{graphicx}
\usepackage{float}%允许浮动
\usepackage{caption}%设置标题
\usepackage{graphicx, subfig}

%附录
\usepackage[toc,page]{appendix}

%导言区
\title{中小微企业的信贷决策}

%不要页眉
\pagestyle{plain}
\bibliographystyle{plain}

\geometry{a4paper}%将纸张设置为A4大小

%修改文档结构层次的名称
\usepackage{titlesec}%chapter1修改为一
\titleformat{\part}{\centering\Huge\bfseries}{\thepart}{1em}{}
\renewcommand\thepart{\chinese{part}}

%subsubsubsection的使用
%\usepackage{titlesec}
%\titleclass{\subsubsubsection}{straight}[\subsection]
%
%\newcounter{subsubsubsection}[subsubsection]
%\renewcommand\thesubsubsubsection{\thesubsubsection.\arabic{subsubsubsection}}
%\renewcommand\theparagraph{\thesubsubsubsection.\arabic{paragraph}}
%
%\titleformat{\subsubsubsection}
%{\normalfont\normalsize\bfseries}{\thesubsubsubsection}{1em}{}
%\titlespacing*{\subsubsubsection}
%{0pt}{3.25ex plus 1ex minus .2ex}{1.5ex plus .2ex}
%
%\makeatletter
%\renewcommand\paragraph{\@startsection{paragraph}{5}{\z@}%
%	{3.25ex \@plus1ex \@minus.2ex}%
%	{-1em}%
%	{\normalfont\normalsize\bfseries}}
%\renewcommand\subparagraph{\@startsection{subparagraph}{6}{\parindent}%
%	{3.25ex \@plus1ex \@minus .2ex}%
%	{-1em}%
%	{\normalfont\normalsize\bfseries}}
%\def\toclevel@subsubsubsection{4}
%\def\toclevel@paragraph{5}
%\def\toclevel@paragraph{6}
%\def\l@subsubsubsection{\@dottedtocline{4}{7em}{4em}}
%\def\l@paragraph{\@dottedtocline{5}{10em}{5em}}
%\def\l@subparagraph{\@dottedtocline{6}{14em}{6em}}
%\makeatother
%
%\setcounter{secnumdepth}{4}

\begin{document}
\else
\fi
%分文件开始


\part{问题分析}
\noindent \textbf{{\large 问题一}}

%开始：问题1的问题分析

对于问题一，在数据处理方面，我们用商业智能软件Power BI对附件1中的123家企业的企业信息、进项发票信息、销项发票信息进行分析，得到反映经营规模的有效销项金额和反映盈利能力的经营利润率这两项反映企业实力的指标、反映供给侧稳定性的正数有效进项发票率和反映需求侧稳定性的正数有效销项发票率这两项反映企业供求关系稳定性的指标，再将企业信誉评级与违约与否进行量化。

根据数据处理的结果，我们运用层次分析法，构建层次分析模型，求出上述各项指标在评测信贷风险中的权重，计算得出123家企业的信贷能力，以反映这123家企业的信贷风险。接着，我们用MATLAB构建规划模型，借助前述所分析出来的信誉评级，按信贷能力排出的企业排位百分比这两项为依据，判定银行是否给企业贷款，我们还将企业的最大贷款数额和年利率根据排位百分比进行分档，一方面控制企业的贷款额度在10万元到100万元之间，另一方面，使得信誉高、信贷风险小的企业得到利率优惠和更高的贷款额度，随后，我们以银行收入的数学期望为目标函数，以每家企业贷款总额之和小于等于银行年度信贷总额为约束条件，使目标函数最大。

我们求解该模型即得到银行在年度信贷总额固定时对企业的信贷策略。结果以表格的形式呈现了不同企业在不同年度信贷总额下的信贷数额。

%结束：问题1的问题分析

~\\

\noindent \textbf{{\large 问题二}}

%开始：问题2的问题分析
对于问题二，我们在问题一的基础上，通过进一步分析附件2中给出的302家企业的进项与销项发票信息，计算出其经营规模、盈利能力、供给侧稳定性及需求稳定性的数据指标。接着，利用MPai数据科学平台，配置决策树模型，根据已有信誉评级设置相应工程参数，在分析结果页查看最终训练结果，快速得出最拟合数据的模型特征概要，推断出302家企业对应的信誉评级。之后，同问题一的思路，我们运用层次分析理论，对信贷风险进行量化分析。最后，借助规划求解理论构建信贷策略模型并且在MATLAB中实现，得出该银行在年度信贷总额为1亿元时对这些企业的信贷策略。

%结束：问题2的问题分析

~\\

\noindent \textbf{{\large 问题三}}

%开始：问题3的问题分析
对于问题三，在信息处理方面，为了分析新冠病毒疫情对不同行业不同类别的企业的经济效益的影响，我们将302家企业根据行业类别进行分类，得到个体经营、服务业、信息传输业、医疗业等13个行业，并且在调查了大量研究和新闻的基础上，根据行业类别判定受疫情影响的程度，分为有利、基本无影响、不利、非常不利4类。

根据信息处理的结果，我们在问题二的基础上，通过上述影响企业实力、信誉的指标及受疫情影响的程度对企业的信贷风险进行量化分析。最后，借助规划求解理论在MATLAB中编写代码并调用相关函数，构建信贷策略模型得出在疫情影响下该银行在年度信贷总额为1亿元时的信贷调整策略。
%结束：问题3的问题分析



%分文件结束
\ifx\allfiles\undefined
\end{document}
\fi